% !PRINTED 397
\weberpage{397}
%<<CONT
\[
Q_{n-1} \leqq Q_n,
\]
where the equality sign is possible only in the case of the lower
sign, that is, in the case of equation (9).

Further (since $b < \sqrt{D}$) we have
\[
Q_{n-1} = \frac{t-ub}{2} > \frac{t-u\sqrt{D}}{2} = \frac{t^2-u^2D}{2(t+u\sqrt{D})} = \frac{\mp 2}{t+u\sqrt{D}},
\]
so $Q_{n-1}$ is zero or positive, and zero can occur only in the case
of the upper sign, that is, in the case of equation (8).
From this we conclude
\weq{12}{0 \leqq Q_{n-1} \leqq Q_n,}
where the equality sign can occur at the lower limit only in the
case of equation (8), that is, of the upper signs, and at the upper
limit only in the case of equation (9), that is, of the lower signs.

We now develop the rational fraction $P_n : Q_n$ into a finite
continued fraction
\[
\frac{P_n}{Q_n} = (\alpha_0, \alpha_1, \alpha_2 \ldots \alpha_{n-1}),
\]
by taking $n$, according to \S.\,115, I, so that
\[
(-1)^n = \mp 1
\]
If then $P' : Q'$ denotes the next-to-last convergent, we have
\[
P_n Q' - Q_n P' = \mp 1
\]
and, by \S.\,117, II:
\[
0 \leqq Q' \leqq Q_n,
\]
where the equality sign can occur at the lower limit only for $n=1$,
that is, in the case of the upper signs, and at the upper limit only
for $n=2$, that is, of the lower signs, just as in (12). From this
it follows by \S.\,118, III:
\[
Q' = Q_{n-1}, \quad P' = P_{n-1},
\]
and hence, by (10),
\[
(\alpha_0, \alpha_1, \alpha_2 \ldots \alpha_{n-1}, \omega) = \omega,
\]
that is, $[\alpha_0, \alpha_1, \alpha_2 \ldots \alpha_{n-1}]$ is a period of the continued fraction for
$\omega$, from which, by the formulae (4), (5), the solution $t, u$ of (7)
is derived. It is thereby proved that in this way we find
\gesperrt{all solutions} of equation (7).

From this some further important consequences can be drawn.
